Despite its realism and scale, this work has several limitations. First, RetailBench focuses on a single-store supermarket setting; while expressive, it does not capture multi-store coordination, competitive markets, or strategic interactions among multiple autonomous agents. Second, although the environment incorporates stochastic demand, news dynamics, and supply-chain delays, it remains a simulation grounded in historical data and simplified economic assumptions, which may not fully reflect the complexities of real-world retail systems.

Third, our evaluation is limited to prompting-based LLM agents without parameter updates or long-term learning across episodes; stronger performance may be achievable through reinforcement learning, fine-tuning, or hybrid neuro-symbolic approaches. Finally, while we identify key failure modes such as hallucinations and economically irrational actions, we do not propose explicit algorithmic mechanisms to enforce economic constraints or factual grounding during execution. Addressing these limitations—through richer environments, multi-agent extensions, learning-based adaptation, and constraint-aware action control—remains an important direction for future research.